\chapter{\texorpdfstring{The Quantum Balance (Brazil):
l'\emph{Estado Logístico} come
\(\Phi\)-engine}{The Quantum Balance (Brazil): l'Estado Logístico come \textbackslash Phi-engine}}\label{il-bilanciere-quantistico-brasile-lestado-loguxedstico-come-phi-engine}

\begin{quote}
\textbf{Core argument:} il Brazil del 2026 è l'unico fra i quattro
attori della Part III a possedere la \emph{struttura formal} di un
produttore di \(\Phi\) esportabile --- non per dotazione materiale
(\(s_{ii}\) medio, non superiore), né per egemonia ideologica
(\(s_{iii}\) regionale-globale ma non globale-totale), ma per
\textbf{assetto operational}: l'\emph{Estado Logístico} di Amado Cervo
come politica estera che mantiene tutti i fronti aperti, la capacità di
mediazione plurale documentata su Russia-Ucraina, BRICS, COP30, e una
resilience democratica istituzionale verificationta (V-Dem U-turn 2025) after
aver respinto il colpo di stato del gennaio 2023 e condannato il suo
principale agente. La \(\Phi\)-engine Brazilian è however
\textbf{conditional}: l'elezione presidenziale del 4 ottobre 2026 è il
test di falsifiability più stretto della DEQ$^2$ nell'orizzonte 2028-2031.
\end{quote}



\section{Structural Definition del Bilanciere
Quantistico}\label{definition-structural-del-bilanciere-quantistico}

\Hypoth{} \textbf{Definizione (Quantum Balance DEQ$^2$).} Un attore
\(e\) è detto \emph{Quantum Balance} (o \(\Phi\)-engine) se
soddisfa simultaneamente:

\begin{enumerate}
\def\labelenumi{\arabic{enumi}.}
\tightlist
\item
  \(\Phi_e^{\text{interno}}\) moderato ma \textbf{prodotto
  spontaneamente} (non per coercizione gerarchica come nel Rigido
  Sincronizzato, non assente come nell'Decoherent Hegemon);
\item
  \(\Phi_e^{\text{esportabile}} > 0\) --- esistono \emph{interfacce di
  coupling bidirezionale} con sotto-sistemi globali eterogenei
  (USA, China, Russia, UE, Africa, Medio Oriente, Asia) che permettono di
  trasferire ordine senza imporre allineamento;
\item
  \textbf{Asimmetria \(A\)/\(\Omega\) favorevole}: \(A_e\) adattivo
  (esperienza ricorrente di crisi superate senza collapse
  istituzionale), \(\Omega_e\) sotto-medio per la fascia (l'energia
  dissipata in conflitto interno cresce ma resta entro la soglia di
  gestibilità);
\item
  \(s_{iii,e} > s_{ii,e}\) in termini relativi (la legitimacy
  etico-discorsiva eccede la material capacity --- caratteristica
  structural di un \emph{bridge-builder});
\item
  \textbf{Politica estera non-allineata operativa}: l'Estado Logístico
  di Cervo come \emph{policy} effettiva, non come retorica.
\end{enumerate}

\Proved{} \textbf{Confollowsnza formal.} \(T_s^{(2),\text{sys}}\) del
Bilanciere è strutturalmente in regime \textbf{stabile-perturbabile}:
alta sotto condizioni di base, sensibile a tre variables --- esito
elettorale, fiscale, polarizzazione interna --- il cui collapse
simultaneo produrrebbe una transizione verso \emph{Decoherent Hegemon}
(caso USA-like) o \emph{Soliton Politico} (caso post-2018). Non esiste
rotta verso \emph{Rigid Synchronized}, because lo \emph{stack}
istituzionale federale Brazilian è asimmetrico with respect to quello
Chinese.



\section{\texorpdfstring{L'\emph{Estado Logístico} di Amado Cervo come
Apparato
Operativo}{L'Estado Logístico di Amado Cervo come Apparato Operativo}}\label{lestado-loguxedstico-di-amado-cervo-come-apparato-operational}

\Hypoth{} \textbf{Eredità teorica.} L'\emph{Estado Logístico} (Amado
Luiz Cervo, \emph{Inserção internacional: formação dos conceitos
brasileiros}, 2008) è il paradigma di politica estera in cui lo Stato
non rinuncia all'autonomia decisionale ma la articola through una
\emph{molteplicità simultanea di vettori}: industriale (BRICS), agricolo
(Cerrado, soia), energetico (idroelettrico, biocombustibili, oggi
rinnovabili), diplomatico (mediazione plurale), ambientale (Amazzonia,
COP). Cervo lo distingue da \emph{Estado Normal} (subordinato),
\emph{Estado Desenvolvimentista} (autarchico), \emph{Estado Neoliberal}
(dipending).

\Proved{} \textbf{Mappatura DEQ$^2$ formal.} L'\emph{Estado Logístico} è
la realizzazione politica della condizione (5) della definition
\$\S\$9.0. In termini del campo \(\Pi\) topografico:

\[\Pi^{\text{Logístico}}(t+1) = \sum_{k} M_e^{(k)} \cdot Q_0^{(k)} \cdot (1 - \delta_R^{(k)} z_R^{(k)})\]

where \(k\) indicizza i diversi vettori (industriale, agricolo,
energetico, diplomatico, ambientale) e ciascun \(M_e^{(k)}\),
\(Q_0^{(k)}\), \(\delta_R^{(k)}\) è specifico del vettore. È una
\emph{sovrapposizione coerente} --- esattamente la struttura matematica
che genera \(\Phi\) in un sistema multi-canale.

\Hypoth{} \textbf{Confollowsnza structural.} Il Bilanciere produce
\(\Phi\) globale per \emph{sovrapposizione plurale} (più vettori
coerenti) e non per \emph{amplificazione singolare} (un solo vettore
massimizzato). Questo lo rende qualitativamente diverso dal Soliton
(Ch.~6), dall'Decoherent Hegemon (Ch.~7) e dal Rigid Synchronized
(Ch.~8). Strutturalmente, il Bilanciere è \textbf{l'unico attore della
Part III in cui l'apparato matematico DEQ$^2$ produce numeratore \(> 0\)
stabile}.



\section{BRICS 2025 Rio + COP30 Belém: due Test
Simultanei}\label{brics-2025-rio-cop30-beluxe9m-due-test-simultanei}

\Proved{} \textbf{Dato 1 --- Presidenza BRICS 2025.} Brazil presiede
BRICS dal 1 gennaio 2025. Il 17$^\circ$ summit (Rio de Janeiro, 6-7 luglio
2025) approva tre documenti centrali della presidenza Brazilian: -
\emph{BRICS Leaders' Framework Declaration on Climate Finance}; -
\emph{BRICS Leaders' Declaration on Global Governance of Artificial
Intelligence}; - \emph{BRICS Partnership for the Elimination of Socially
Determined Diseases}.

Dichiarazione finale: 126 punti su global governance, finanza, salute,
AI, clima. Lancio del \emph{Tropical Forest Forever Fund} (TFFF) e del
\emph{BRICS Multilateral Guarantees} (BMG).

\Proved{} \textbf{Dato 2 --- COP30 Belém.} 10-22 novembre 2025. Brazil
presiede fino a novembre 2026. Tre pilastri: 1. \emph{Strengthening
multilateralism and the climate regime}; 2. \emph{Connecting climate
initiatives to people's daily lives}; 3. \emph{Accelerating
implementation of the Paris Agreement}.

Risultati: mobilitazione \(1{,}3\)T USD/anno entro 2035, triplicamento
finanza adattamento, \emph{Belém Mechanism for Just Global Transition},
due roadmap (transizione fuori dai fossili + foreste/clima) lanciate per
iniziativa della presidenza Brazilian malgrado assenza di consenso
(\(\approx 80\) paesi pro / \(\approx 80\) contro language esplicito).

\Hypoth{} \textbf{Lettura DEQ$^2$ congiunta.} I due eventi, ravvicinati
(luglio 2025 e novembre 2025), misurano la \(\Phi\) esportabile
Brazilian in due assi distinti:

{\def\LTcaptype{none} % do not increment counter
{\small\begin{longtable}[]{@{}
  >{\raggedright\arraybackslash}p{(\linewidth - 4\tabcolsep) * \real{0.3333}}
  >{\raggedright\arraybackslash}p{(\linewidth - 4\tabcolsep) * \real{0.3333}}
  >{\raggedright\arraybackslash}p{(\linewidth - 4\tabcolsep) * \real{0.3333}}@{}}
\toprule\noalign{}
\begin{minipage}[b]{\linewidth}\raggedright
Test
\end{minipage} & \begin{minipage}[b]{\linewidth}\raggedright
Asse misurato
\end{minipage} & \begin{minipage}[b]{\linewidth}\raggedright
Risultato 2025
\end{minipage} \\
\midrule\noalign{}
\endhead
\bottomrule\noalign{}
\endlastfoot
BRICS Rio & \(\Phi^{\text{Sud globale}}\) --- capacità di sintesi
multilaterale fra interessi divergenti (China-India, Russia-Iran,
Brazil-mediatore) & Documenti tematici di consenso; iniziative
innovative (TFFF, BMG); 200+ riunioni preparatorie \\
COP30 Belém & \(\Phi^{\text{Nord-Sud}}\) --- capacità di mediare fra
alta-emissioni e adattamento, mantenendo entrambi nel campo \(\Pi\) &
Pacchetto Belém approvato; due roadmap lanciate \emph{malgrado} assenza
di consenso (mossa di leadership) \\
\end{longtable}}
}

\Proved{} \textbf{Confollowsnza.} Il Brazil è l'unico Stato del 2025 ad
aver presieduto contemporaneamente un blocco multilaterale Sud (BRICS) e
un consesso multilaterale globale (COP), producendo in entrambi
documenti operativi non banali. Questa è la firma empirica
dell'asimmetria definita in \$\S\$9.0 condizione 4: \(s_{iii} > s_{ii}\)
produce \emph{capacità di convocazione} superiore alla \emph{capacità di
imposizione}.

\Conj{} \textbf{Caveat onesto.} L'efficacia dei documenti del 2025 si
misura nei trienni successivi. Il TFFF ha bisogno di donazioni
effettive; il BMG di investitori reali; le roadmap COP di adesioni. Il
Brazil ha \emph{aperto i fronti} --- l'esito misurerà se il vettore
\(\Phi\)-esportabile è solo dichiarativo o effettivamente operational.



\section{Capacità di Mediazione come Struttura, non
Episodio}\label{capacituxe0-di-mediazione-come-struttura-non-episodio}

\Proved{} \textbf{Dato.} Maggio 2025: Lula visita Mosca, propone
\emph{strategic partnership} RU-BR e si offre come mediatore RU-UA,
condizionatamente all'apertura di entrambe le parti. Il Brazil è
co-fondatore (con la China) del \emph{Group of Friends for Peace} (gruppo
amici della pace), che ha pubblicato un piano in sei punti per il
conflitto ucraino (2024).

\Hypoth{} \textbf{Lettura DEQ$^2$ della mediazione structural.} La
capacità di mediazione del Brazil non è un attributo personale di Lula
ma una \emph{invariante structural} dell'Estado Logístico: dipende dal
mantenimento attivo di tutti i canali di comunicazione (USA, China,
Russia, UE) senza chiusura ideologica unilaterale. Formalmente: il
Brazil mantiene il proprio campo \(\Pi^{\text{globale}}\) con
\(\Sigma\) (matrice di covarianza degli assi \(x\) coerenza, \(y\)
precisione, \(z\) dialectical complexity --- eredità Topografia \S{}5.1)
\textbf{massimamente aperta}, contro l'Decoherent Hegemon che la riduce
a \(\Sigma^{\text{egemonica}}\) con \(z \to 0\) (cap. 7.3).

\Proved{} \textbf{Confollowsnza measurable.} Il Brazil è l'unico fra i
grandi Stati a: - non aver applicato sanzioni a Russia, China, Iran; -
mantenere relazioni operative con Israele e Iran simultaneamente; -
preservare partnership economiche con USA e China simultaneamente; -
partecipare a BRICS+ senza abbandonare OECD-track e G20.

\Hypoth{} \textbf{Tesi formal.} Il \emph{non-allineamento attivo} è la
condizione operativa della \(\Phi\)-engine: produrre coerenza globale
senza imporre allineamento sub-globale. È una posizione sistemicamente
distinta dal \emph{non-allineamento passivo} della Guerra Fredda (India,
Yugoslavia) --- è \emph{non-allineamento generativo}, capace di produrre
nuovi accoppiamenti (TFFF, \emph{Group of Friends for Peace},
\emph{Belém Mechanism}).



\section{\texorpdfstring{V-Dem U-Turn: la Prova di \(\Phi\)
Resiliente}{V-Dem U-Turn: la Prova di \textbackslash Phi Resiliente}}\label{v-dem-u-turn-la-prova-di-phi-resilient}

\Proved{} \textbf{Dato V-Dem 2026} (relativo all'anno 2025): registrato
un \emph{U-turn} (inversione di autocratizzazione) per il Brazil
post-2022. La presidenza Lula ha \textbf{invertito il deterioramento}
documentato sotto Bolsonaro. Il rapporto annote: miglioramenti
istituzionali, \emph{judicial assertiveness} (STF, TSE), \emph{civil
society mobilization}, \emph{media independence}, \emph{neutrality of
armed forces}. Avvertimento esplicito: la \emph{polarizzazione sociale
persiste} e \emph{l'elezione del 2026 sarà decisiva}.

\Proved{} \textbf{Verifica empirica della resilience istituzionale.}
L'11 settembre 2025 la Suprema Corte Federale (STF) condanna Jair
Bolsonaro a 27 anni e 3 mesi per: - tentativo di colpo di stato (8
gennaio 2023); - partecipazione a organizzazione criminale armata; -
piano per uccidere Lula, Alckmin, Moraes; - abolizione violenta
dell'ordine democratico.

7 novembre 2025: pannello STF respinge l'appello (4-1, dissenso del solo
giudice Fux). È la \textbf{first condanna nella storia Brazilian di un
ex-presidente per tentato colpo di stato}.

\Hypoth{} \textbf{Lettura DEQ$^2$. } La condanna Bolsonaro è l'evento
empirico più rilevante della Part III per la teoria DEQ$^2$: misura la
capacità del sistema Brazilian di mantenere
\(\Phi^{\text{interno}} > 0\) sotto perturbazione massima (tentativo di
rottura dell'ordine democratico). In termini formali, ha mostrato che il
sistema possiede un \emph{meccanismo di feedback negativo} funzionante
--- una proprietà structural che né Decoherent Hegemon né Rigido
Sincronizzato esibiscono.

{\def\LTcaptype{none} % do not increment counter
{\small\begin{longtable}[]{@{}
  >{\raggedright\arraybackslash}p{(\linewidth - 6\tabcolsep) * \real{0.2500}}
  >{\raggedright\arraybackslash}p{(\linewidth - 6\tabcolsep) * \real{0.2500}}
  >{\raggedright\arraybackslash}p{(\linewidth - 6\tabcolsep) * \real{0.2500}}
  >{\raggedright\arraybackslash}p{(\linewidth - 6\tabcolsep) * \real{0.2500}}@{}}
\toprule\noalign{}
\begin{minipage}[b]{\linewidth}\raggedright
Indicatore istituzionale
\end{minipage} & \begin{minipage}[b]{\linewidth}\raggedright
USA 2026
\end{minipage} & \begin{minipage}[b]{\linewidth}\raggedright
China 2026
\end{minipage} & \begin{minipage}[b]{\linewidth}\raggedright
Brazil 2026
\end{minipage} \\
\midrule\noalign{}
\endhead
\bottomrule\noalign{}
\endlastfoot
Vincoli giudiziari sull'esecutivo & \(\downarrow\) minimo da 100+ anni
(Ch.~7.1) & non applicabili & \(\uparrow\) STF assertiva \\
Indipendenza media & \(\downarrow\) (attacchi a stampa) & bassa per
costruzione & \(\uparrow\) post-Bolsonaro \\
Mobilitazione civile & \(\downarrow\) frammentata & \(\to 0\) permessa &
media-alta \\
Neutralità forze armate & erosa nel 2021-2024 & partito = forze armate &
confermata 2023-2025 \\
\end{longtable}}
}

\Proved{} \textbf{Confollowsnza.} Il Brazil possiede ciò che la
Topografia \S{}5.4 chiama \emph{critical resistance} \(R_c\): un sistema di
sotto-attori che possono opporsi alla traiettoria di un singolo attore
dominante. Per la DEQ$^2$, \(R_c > 0\) effettivo è condizione necessaria di
\(\Phi\) spontanea (vs \(\Phi\) imposta del Ch.~8).



\section{Le Fragilità: Onestà
Analitica}\label{le-fragilituxe0-onestuxe0-analitica}

Il Ch.~9 perderebbe rigore se omettesse le fragility reali del
Bilanciere Brazilian. La \(\Phi\)-engine non è automatica --- è
\textbf{conditional} a tre variables in tensione.

\subsection{Fragilità 1: Polarizzazione e approval
calante}\label{fragilituxe0-1-polarizzazione-e-approval-calante}

\Proved{} \textbf{Dato.} Approval Lula marzo 2026 (Genial/Quaest, 2.004
elettori): \(44\%\) approva, \(51\%\) disapprova --- peggior gap dal
luglio 2025. Sondaggi alternativi danno \(28\%/40\%\). Coalizione di
governo a \emph{minimo storico di 30 anni} in fedeltà parlamentare.

\Conj{} \textbf{Lettura DEQ$^2$.} La polarizzazione interna è il rumore
endogeno che erode \(\Phi^{\text{interno}}\). Il Brazil non è immune al
fenomeno globale di affective polarization --- è solo che dispone di
meccanismi \(R_c\) ancora operanti per contenerlo.

\subsection{Fragilità 2: Fiscale e
Selic}\label{fragilituxe0-2-fiscale-e-selic}

\Proved{} \textbf{Dato.} Crescita 2025: \(+2{,}3\%\). Forecast 2026:
\(+1{,}7\%\). Selic da \(10{,}5\%\) (settembre 2024) a \(15\%\) (giugno
2025); easing atteso fine 2026 attorno a \(9{,}75\)-\(11{,}50\%\).
Inflazione IPCA 2025: \(4{,}44\%\) (entro target \(1{,}5\)-\(4{,}5\%\)).
Deficit federale 2026: \(\approx 3\%\) PIL; debito pubblico
\(\approx 75\%\) PIL. Quadro fiscale \emph{weakened} da indebolimento
delle regole proprie.

\Conj{} \textbf{Lettura DEQ$^2$.} Lo squilibrio fiscale aumenta \(\sigma\)
esogena --- volatilità monetaria, pressione su tasso di cambio,
vulnerabilità a shock esterni. Il Brazil resta sotto la soglia di
\emph{currency crisis}, ma la traiettoria 2025-2026 è sub-ottimale per
un \(\Phi\)-engine. Strutturalmente: la combinazione \emph{crescita
rallentata + Selic alto + debito 75\%} riduce la capacità del governo di
sostenere finanziariamente l'agenda Estado Logístico (TFFF, mediazione,
COP).

\subsection{Fragilità 3: Elezioni 4 ottobre
2026}\label{fragilituxe0-3-elezioni-4-ottobre-2026}

\Proved{} \textbf{Dato.} Bolsonaro ineligibile (condannato + in
prigione). Suo figlio \textbf{Flávio Bolsonaro} candidato del Partido
Liberal. Sondaggio AtlasIntel/Bloomberg aprile 2026: Flávio \(47{,}6\%\)
vs Lula \(46{,}6\%\) in secondo turno simulato --- \textbf{first volta}
in cui Flávio è numericamente avanti (margine entro errore statistico).
Primo turno: Lula \(45{,}9\%\) vs Flávio \(40{,}1\%\). Trend dei mesi
precedenti: Lula in calo da \(+12\) pp (dicembre 2025) a pareggio
statistico.

\Conj{} \textbf{Lettura DEQ$^2$ centrale.} L'elezione del 4 ottobre 2026 è
il \textbf{test di falsification più stretto} della DEQ$^2$ nell'orizzonte
2028-2031. Tre scenari:

{\def\LTcaptype{none} % do not increment counter
{\small\begin{longtable}[]{@{}
  >{\raggedright\arraybackslash}p{(\linewidth - 4\tabcolsep) * \real{0.3333}}
  >{\raggedright\arraybackslash}p{(\linewidth - 4\tabcolsep) * \real{0.3333}}
  >{\raggedright\arraybackslash}p{(\linewidth - 4\tabcolsep) * \real{0.3333}}@{}}
\toprule\noalign{}
\begin{minipage}[b]{\linewidth}\raggedright
Scenario
\end{minipage} & \begin{minipage}[b]{\linewidth}\raggedright
Probabilità soggettiva (aprile 2026)
\end{minipage} & \begin{minipage}[b]{\linewidth}\raggedright
Confollowsnza per \(\Phi\)-engine
\end{minipage} \\
\midrule\noalign{}
\endhead
\bottomrule\noalign{}
\endlastfoot
\textbf{A} Lula re-eletto & \(0{,}45\)-\(0{,}55\) & \(\Phi\)-engine
continua, sotto pressione fiscale ma operativa \\
\textbf{B} Successore di sinistra (post-Lula, es. Haddad) &
\(0{,}10\)-\(0{,}15\) & \(\Phi\)-engine continua, possibile riforma
fiscale \\
\textbf{C} Flávio Bolsonaro o destra radicale eletta &
\(0{,}30\)-\(0{,}40\) & \(\Phi\)-engine si interrompe, traiettoria verso
\emph{Decoherent Hegemon regionale} \\
\end{longtable}}
}

\Hypoth{} \textbf{Tesi forte.} Se A o B: la DEQ$^2$ prevede consolidamento
del Brazil come \(\Phi\)-engine globale entro 2028-2031, con
\(T_s^{(2),\text{sys}}\) nettamente superiore a USA e China. Se C: la
DEQ$^2$ prevede solitonizzazione discorsiva del Brazil (passaggio dal
profilo Dialogico \((7,6,8)\) all'Egemonico \((8,9,2)\) --- ricalcando
la traiettoria USA del Ch.~7.3) con perdita parziale di \(s_{iii}\)
globale entro 24 mesi.



\section{\texorpdfstring{Chapter Summary: \(\Phi\)-engine
Condizionale}{Chapter Summary: \textbackslash Phi-engine Condizionale}}\label{sintesi-del-chapter-phi-engine-conditional}

\Proved{} \textbf{In una frase:} il Brazil del 2026 è strutturalmente
l'unico attore della Part III configureto come \(\Phi\)-engine ---
Estado Logístico operational, doppia presidenza BRICS-COP, mediazione
structural, V-Dem U-turn verified, condanna Bolsonaro come prova
empirica di \(R_c\) effettivo --- ma la condizionalità elettorale del 4
ottobre 2026 fa della DEQ$^2$ una teoria \emph{con punto di bifurcation
esplicito}: il sistema globale entrerà nella finestra 2028-2031 con o
senza un produttore di \(\Phi\) esportabile.

{\def\LTcaptype{none} % do not increment counter
{\small\begin{longtable}[]{@{}
  >{\raggedright\arraybackslash}p{(\linewidth - 6\tabcolsep) * \real{0.2500}}
  >{\raggedright\arraybackslash}p{(\linewidth - 6\tabcolsep) * \real{0.2500}}
  >{\raggedright\arraybackslash}p{(\linewidth - 6\tabcolsep) * \real{0.2500}}
  >{\raggedright\arraybackslash}p{(\linewidth - 6\tabcolsep) * \real{0.2500}}@{}}
\toprule\noalign{}
\begin{minipage}[b]{\linewidth}\raggedright
Variabile
\end{minipage} & \begin{minipage}[b]{\linewidth}\raggedright
Valore stimato 2026
\end{minipage} & \begin{minipage}[b]{\linewidth}\raggedright
Direzione
\end{minipage} & \begin{minipage}[b]{\linewidth}\raggedright
Confollowsnza per \(T_s^{(2),\text{sys}}\)
\end{minipage} \\
\midrule\noalign{}
\endhead
\bottomrule\noalign{}
\endlastfoot
\(s_{ii}\) (material capacity) & media & stabile & numeratore
moderato \\
\(s_{iii}\) (ethical-discursive legitimacy) & \(49{,}2\) Brand Finance,
V-Dem U-turn & \(\uparrow\) & numeratore in crescita \\
\(\Phi^{\text{interno}}\) & spontaneo, sotto pressione polarizzazione &
stabile-perturbabile & numeratore funzionale \\
\(\Phi^{\text{esportabile}}\) & \(> 0\) verified (BRICS, COP30,
mediazione) & \(\uparrow\) & unicità nel sistema \\
\(A^{\text{adattivo}}\) & alto (resilience 2018-2023 confermata) &
\(\uparrow\) & denominatore contenuto \\
\(\Omega\) (dissipazione interna) & medio (polarizzazione persistente) &
stabile & denominatore moderato \\
\(\varepsilon_{\text{dis}}\) (disimpegno morale) & calo post-2022,
possibile risalita 2026 & dipende elezioni & swing \\
\textbf{Variabile di bifurcation} & esito 4 ottobre 2026 &
\(P(\text{Lula/sinistra}) \approx 0{,}55\)-\(0{,}70\) & A/B vs C separa
due regimi \\
\end{longtable}}
}

\Hypoth{} \textbf{Tesi conclusiva del chapter.} La probabilità che il
sistema globale 2028-2031 entri nella sua phase transition con almeno
un produttore di \(\Phi\) esportabile è \(P \approx 0{,}55\)-\(0{,}70\),
conditional al risultato elettorale Brazilian dell'ottobre 2026. È la
singola variable più informativa fra tutte quelle considerate nei Cap.
6-9 --- because tutti gli altri attori (Soliton, USA, China) hanno già
\emph{firmato} la propria traiettoria DEQ$^2$, while il Brazil non l'ha
ancora fatto.

\Hypoth{} \textbf{Asimmetria triplice consolidata (Ch.~7-8-9).}

{\def\LTcaptype{none} % do not increment counter
{\small\begin{longtable}[]{@{}
  >{\raggedright\arraybackslash}p{(\linewidth - 6\tabcolsep) * \real{0.2500}}
  >{\raggedright\arraybackslash}p{(\linewidth - 6\tabcolsep) * \real{0.2500}}
  >{\raggedright\arraybackslash}p{(\linewidth - 6\tabcolsep) * \real{0.2500}}
  >{\raggedright\arraybackslash}p{(\linewidth - 6\tabcolsep) * \real{0.2500}}@{}}
\toprule\noalign{}
\begin{minipage}[b]{\linewidth}\raggedright
Attore
\end{minipage} & \begin{minipage}[b]{\linewidth}\raggedright
Tipo \(\Phi\)
\end{minipage} & \begin{minipage}[b]{\linewidth}\raggedright
Esportabilità
\end{minipage} & \begin{minipage}[b]{\linewidth}\raggedright
Direzione \(T_s^{(2),\text{sys}}\)
\end{minipage} \\
\midrule\noalign{}
\endhead
\bottomrule\noalign{}
\endlastfoot
Soliton (Musk) & \(\Phi\) degenere (\(N=1\)) & per decoherence forzata
su altri & indefinita \\
Decoherent Hegemon (USA) & \(\Phi\) mai prodotta & nessuna & \(\to 0\)
in 2027-2029 \\
Rigid Synchronized (China) & \(\Phi\) imposta & esportazione di \(K\),
non di \(\Phi\) & \(\to 0\) catastrofica \\
\textbf{Quantum Balance (Brazil)} & \(\Phi\) spontanea
conditional & \textbf{vera \(\Phi\)-engine se A/B vince} &
\textbf{\(> 0\) se A/B; \(\to 0\) se C} \\
\end{longtable}}
}



\section{Quattro Predizioni Datate del Ch.~9 (per Cap.
11)}\label{quattro-predizioni-datate-del-cap.-9-per-cap.-11}

\begin{enumerate}
\def\labelenumi{\arabic{enumi}.}
\tightlist
\item
  \textbf{Esito elezione 4 ottobre 2026 condiziona la traiettoria.}
  Predizione pubblica:
  \(P(\text{vittoria Lula o successore di sinistra}) = 0{,}55\)-\(0{,}70\)
  aprile 2026, da rivedere al 30 giugno 2026 e al 30 settembre 2026
  sulla base di sondaggi aggregati (Genial/Quaest, AtlasIntel,
  Datafolha).
\item
  \textbf{Se Lula/sinistra vince}: TFFF raccoglie \(> 5\)B USD impegnati
  entro fine 2027 e BRICS Multilateral Guarantees parte operativamente
  entro Q2 2027. Falsificabile su comunicati BRICS e World Bank.
\item
  \textbf{V-Dem Brazil in salita}: indice Liberal Democracy oltre
  \(0{,}65\) entro report 2028 (su anno 2027), conditional alla
  continuità Lula/successore. Falsificabile su \url{https://v-dem.net}.
\item
  \textbf{Se vince Flávio Bolsonaro o destra radicale}: deterioramento
  V-Dem \(> 0{,}05\) punti in 24 mesi; Brazil esce da \emph{Group of
  Friends for Peace}; partecipazione attiva BRICS si trasforma in
  adesione formal ma passiva. Falsificabile su corso degli eventi
  2027-2028.
\end{enumerate}



\section{Closing Epistemic Notes
Chapter}\label{note-epistemiche-di-chiusura-chapter}

{\def\LTcaptype{none} % do not increment counter
{\small\begin{longtable}[]{@{}
  >{\raggedright\arraybackslash}p{(\linewidth - 4\tabcolsep) * \real{0.3333}}
  >{\raggedright\arraybackslash}p{(\linewidth - 4\tabcolsep) * \real{0.3333}}
  >{\raggedright\arraybackslash}p{(\linewidth - 4\tabcolsep) * \real{0.3333}}@{}}
\toprule\noalign{}
\begin{minipage}[b]{\linewidth}\raggedright
\#
\end{minipage} & \begin{minipage}[b]{\linewidth}\raggedright
Claim
\end{minipage} & \begin{minipage}[b]{\linewidth}\raggedright
Status
\end{minipage} \\
\midrule\noalign{}
\endhead
\bottomrule\noalign{}
\endlastfoot
1 & Definizione formal Quantum Balance & \Hypoth{} categoria
DEQ$^2$ nuova \\
2 & Estado Logístico (Cervo) come \(\Phi^{\text{esportabile}}\)
operational & \Hypoth{} mappatura concettuale forte \\
3 & Doppia presidenza BRICS Rio + COP30 Belém 2025 come test simultaneo
& \Proved{} dato pubblicato \\
4 & Mediazione RU-UA via Group of Friends for Peace & \Proved{} dato
pubblicato \\
5 & V-Dem 2026 U-turn per Brazil & \Proved{} dato pubblicato \\
6 & Condanna Bolsonaro 27 anni come prova di \(R_c\) & \Proved{} dato
pubblicato 11/09/2025 \\
7 & Approval Lula in calo, fragility reali & \Proved{} dati
Genial/Quaest, Bloomberg \\
8 & Selic \(15\%\), deficit \(\approx 3\%\) PIL 2026, debito
\(\approx 75\%\) & \Proved{} dati BBVA, OECD \\
9 & Sondaggio AtlasIntel: Flávio \(47{,}6\%\) vs Lula \(46{,}6\%\)
secondo turno & \Proved{} dato aprile 2026 \\
10 & Tre scenari elettorali con probabilità
\(A \approx 0{,}45\)-\(0{,}55\), \(B \approx 0{,}10\)-\(0{,}15\),
\(C \approx 0{,}30\)-\(0{,}40\) & \Conj{} stima soggettiva da sondaggi
aggregati \\
11 & \(\Phi\)-engine conditional all'esito elettorale & \Hypoth{}
risultato structural forte \\
12 & Predizioni 1-4 (TFFF, BRICS BMG, V-Dem 2027) & \Conj{}
falsificabili \\
\end{longtable}}
}



\section{Fonti dei dati 2025-2026
utilizzate}\label{fonti-dei-dati-2025-2026-utilizzate}

\begin{itemize}
\tightlist
\item
  Cervo, A. L., \emph{Inserção internacional: formação dos conceitos
  brasileiros} (2008) --- base teorica Estado Logístico
\item
  BRICS Brasil --- 17$^\circ$ summit Rio 6-7/07/2025, Dichiarazione 126 punti,
  TFFF, BMG
\item
  COP30 Belém 10-22/11/2025 --- UN, UNFCCC, Carbon Brief, World
  Resources Institute, House of Commons Library
\item
  V-Dem 2026 --- U-turn Brazil post-2022, \emph{judicial assertiveness}
\item
  STF Brazil --- condanna Bolsonaro 11/09/2025, conferma appello
  7/11/2025 (Al Jazeera, PBS, Time, CNN)
\item
  Genial/Quaest, AtlasIntel/Bloomberg, Datafolha --- sondaggi 2026
\item
  BBVA Research, OECD, Deloitte --- Brazil Economic Outlook 2026
\item
  gov.br/planalto --- Lula visita Mosca 05/2025, Group of Friends for
  Peace
\item
  Brand Finance Soft Power Index 2026 --- Brazil \(49{,}2\)
\item
  Munich Security Report 2025 --- Ch.~8 \emph{Lula Land}
\end{itemize}



\section{Chiusura della Part III e Opening of Part
IV}\label{chiusura-della-parte-iii-e-apertura-della-parte-iv}

I quattro capitoli applicativi della Part III (Soliton, Egemone
Decoerente, Rigid Synchronized, Bilanciere) hanno mappato il sistema
globale 2026 sull'apparato DEQ$^2$. Tre dei quattro attori hanno
traiettorie già definite --- il Soliton verso una metastabilità
esauribile entro 2031, l'Egemone verso \(\Phi \to 0\) distribuito entro
2027-2029, il Rigido verso \(\Phi\) catastrofica entro l'orizzonte
demografico 2046 ma con punti di rottura possibili nel triennio
2027-2029. Il quarto --- il Bilanciere Brazilian --- è in
\emph{bifurcation esplicita} con data e meccanismo noti: 4 ottobre
2026.

La \textbf{Part IV} del libro tratterà la phase transition 2028-2031
considerando le quattro traiettorie come funzioni del tempo accoppiate,
con dashboard SPC-DEQ (Ch.~10) e predizioni datate falsificabili (Cap.
11).


